\documentclass[11pt]{article}
\usepackage[margin=1in]{geometry}
\usepackage{amsmath,amssymb,amsthm,booktabs,hyperref}
\newtheorem{proposition}{Proposition}
\newcommand{\E}{\mathbb E}
\newcommand{\F}{\mathbb F}
\newcommand{\Prb}{\mathbb P}
\title{Decoding the full field values of elliptic determinant moments}
\author{Project research note}
\date{20 November 2026}
\begin{document}
\maketitle

\paragraph{Result.}
The unit-weight moment values can distinguish hits on small curves even
when the number of tuples exceeds the number of available moments.
Six field moments recover the exact witness count for one signed
abscissa pair in each of three slots by a fixed formula.
Planting gives an efficient way to train or validate a target-specific
decoder using fresh inputs with unknown labels. We give a sufficient
condition in terms of a computable partition of the moment values.
We have no partition or other readout proved to meet that condition at
the tuned growing arity. The end-to-end exponent has not improved.

\section{Which distribution must be decoded?}
Let $G=E(\F_p)$, $N=hr=|G|$ on the project's supplied coprime-cofactor
class, and fix a target $T$ in the odd-prime-order subgroup with $2T\ne O$.
Let $b=\lceil\log_2 p\rceil$, $L=\log_2 b$, $g=\log_2 L$ and $M=2^b$.
In slot $i$ a canonical affine abscissa $x$ is retained when
$(\alpha_i x+\beta_i\bmod M)<\tau_i$, with $\alpha_i$ uniform odd,
$\beta_i$ uniform and slots independent. Both signed points are kept.
Write $P$ for the distribution of the full seeds,
$F_i$ for slot $i$, $D=\prod_i|F_i|$, $\rho_i=\tau_i/M$,
$\nu_i=N\rho_i$ and $\lambda=N^{-1}\prod_i\nu_i$.
Use balanced dyadic widths with $1/(2L)<\lambda\le1/L$ and proposed arity
$m\asymp L^{1/4}\sqrt g$. Let $H$ be the ordinary ordered hit count
and $Y=1_{\{H>0\}}$. Uniformly at each prescribed target,
$\E_P H=\lambda(1+o(1))$, $\E_P[H(H-1)]=O(\lambda^2)$ and
$P(Y=1)=\lambda+O(\lambda^2)$ as in
\texttt{../mixed\_coefficient/mixed.tex}.

With $d=m+1$, let $\Delta_{\mathbf a}(T)$ be the $L(dO)$ evaluation
determinant at the tuple and $-T$, in the basis starting with $1,x,\ldots$
from \texttt{target\_incidence.tex}. For regular seeds, the slot
abscissas are disjoint and avoid $x(T)$. Then $\Delta=0$ exactly at a
hit. Irregular seeds have probability $o(\lambda^2)$ asymptotically.
The unit-weight field values in question are
\[
 Z_j(T)=\sum_{\mathbf a\in\prod F_i}\Delta_{\mathbf a}(T)^j\in\F_p.
\]
The previous interpolation gives all orders through $k=\lfloor
2^{b/m^2}\rfloor$ at expected bit cost $2^{(1+o(1))b/m}$.
Include the known integer $D$ as side information when specified.
The results on training below hold for any fixed feature $F$ of these
data. They also hold with independent feature randomness sampled with
the same law on both distributions.

\section{Training at a prescribed target without a fresh label oracle}

Sample $m-1$ independent uniform affine curve points and complete their
sum to $T$ by the unique last point. Retry when it is $O$. The rejected
fraction is at most $1/(N-1)$, by fixing all but one of the sampled
points. Each valid ordered affine tuple of sum $T$ has the same sampling
probability. Uniform affine points can be sampled by rejection on a
uniform abscissa and one of two possible root indices. For each chosen
point, draw the slot multiplier $\alpha_i$ as usual and a uniform
$u_i\in[0,\tau_i)$. Set $\beta_i=u_i-\alpha_i x_i\bmod M$. This gives the
exact seed law conditional on retaining that coordinate.

\begin{proposition}
The resulting marginal seed law $Q$ satisfies $Q(s)=H(s)P(s)/\mu$,
where
\[
 \mu=\E_P H=\frac{(N-1)^m-(-1)^m}{N}\prod_i\rho_i .
\]
For any event $A$ on the observed field values and allowed side information,
the score
\[
 S(A)=\mu+P(A)-2\mu Q(A)
\]
differs from the error of predicting $Y$ by at most $\E_P(H-1)_+$.
For predicting $H\bmod2$ the discrepancy is at most $\E_P[H(H-1)]$.
Both discrepancies are $O(\lambda^2)$.
\end{proposition}

\begin{proof}
Every specified affine tuple survives the independent seeds with the
same probability $\prod_i\rho_i$. Counting its occurrences in the
planting law gives the marginal identity. The number of affine tuples is
the coefficient at nonidentity $T$ in $(\sum_{u\in G}[u]-[O])^m$,
which gives the formula for $\mu$. For either binary truth $B$,
subtract its actual risk from $S(A)$. The difference is
$\E_P[(H-B)(1-2\,1_A)]$. For existence $H-B=(H-1)_+\le H(H-1)/2$.
For parity $0\le H-B\le H(H-1)$.
\end{proof}

For a finite partition of the values into observed bins $v$, population
score minimization selects a bin exactly when $2\mu Q(v)>P(v)$, with
arbitrary ties. In particular $\mu Q(v)/P(v)=\E_P[H\mid v]$ for positive
mass bins. An input with several hits is more likely to be planted.
The score accounts for this bias, including on small examples where
several hits have appreciable probability.

There is an exact asymptotic criterion on the full values themselves.
Write $P_F,Q_F$ for their two laws. Then
\[
 \E_P\operatorname{Var}(H\mid F)
   =\operatorname{Var}_P(H)-\mu^2\chi^2(Q_F\Vert P_F).
\]
This follows by substituting the conditional mean above into the total
variance identity. The least attainable classification error given $F$
is $O(\lambda^2)$ if and only if this conditional variance is
$O(\lambda^2)$. Indeed $\E(H-Y)^2\le\E[H(H-1)]$, and the squared
distance of a binary label to its optimal real regression lies between
half the optimal classification error and that error. Conditional
expectation is an orthogonal projection, so the residual $L^2$
norms for $H$ and $Y$ differ by $O(\lambda)$. Since
$\operatorname{Var}(H)=\mu+O(\lambda^2)$, the equivalent requirement is
$\chi^2(Q_F\Vert P_F)=1/\mu+O(1)$. This measures whether the full
moment information suffices, independently of the cost to decode it.

\paragraph{A sufficient condition within the time budget.}
Suppose a computable compression of the full field values has $s$
possible bins, with $\log_2 s=O(b/m)$, and its optimal existence-bit
error for every prescribed target is $O(\lambda^2)$. Suppose evaluating
the compression and moments together costs $2^{O(b/m)}$ in mean bit
operations under \emph{both} $P$ and $Q$. The previous interpolation
construction has a bound uniform in the particular masks, since scanning
the public widths bounds its slot work. Then one can learn the decision
for every bin at the prescribed target and retain this order of query
cost and error, including training.

In detail, with failure probability at most $\delta$, Hoeffding's
inequality and a union bound over the $2^s$ subsets give uniform
additive $O(\lambda^2)$ accuracy for the score from
\[
 n_P=O\bigl(\lambda^{-4}(s+\log(1/\delta))\bigr),\qquad
 n_Q=O\bigl(\lambda^{-2}(s+\log(1/\delta))\bigr)
\]
independent samples. Use the empirical bin masses, selecting the bins
with negative empirical $P(v)-2\mu Q(v)$. This requires counts and one
pass over the occupied bins. The best empirical score is within
$O(\lambda^2)$ true risk of the best bin decision. Taking
$\delta=O(\lambda^2)$ preserves the averaged guarantee. A prespecified
family of $J$ trial readouts instead costs $\log J$ in the sample counts,
with the time to compare their decisions charged separately. Training
can be repeated per target using private randomness independent of the
test seeds. The sufficient condition demands a bound on approximation
error for the actual hashes. Finding such a compression remains open.

For selection when a good candidate exists, the sample counts can be
reduced. For $J$ fixed candidates put $u=\log(2J/\delta)$ and let
$R(A)$ denote true existence risk. Always $P(A)\le R(A)+O(\lambda)$
and $\mu Q(A^c)\le R(A)+O(\lambda^2)$. Estimate the score as
$\widehat P(A)-\mu+2\mu\widehat Q(A^c)$. Bernstein's inequality for
a Bernoulli frequency $v$ bounds its deviation by
$O(\sqrt{vu/n}+u/n)$ simultaneously over the family. With sufficiently
large constants, $n_P=O(u/\lambda^3)$ and $n_Q=O(u/\lambda)$ therefore
yield $|\widehat S(A)-R(A)|\le R(A)/3+O(\lambda^2)$: use
$2\sqrt{xy}\le c x+y/c$ on the terms involving $R(A)$.
Empirical minimization then has risk at most $2\min_A R(A)+O(\lambda^2)$.
For a prescribed $s$-bin compression one can take $J=2^s$ and still
minimize by bin counts. The earlier larger sample bounds apply when
uniform additive accuracy is desired for every candidate.

\section{An algebraic decoder using six moments}

Take $m=3$ and one abscissa with two signed points in each slot, avoiding
each other's abscissas and $x(T)$. The $L(4O)$ basis is $1,x,x^2,y$.
Expansion along its only ordinate coordinate gives
$\Delta=C+a_1\epsilon_1+a_2\epsilon_2+a_3\epsilon_3$, where the
eight tuples use every sign pattern once. The first six unit-weight
field moments alone determine the product of all eight determinants
and the number that vanish. All the following operations have fixed
cost in field arithmetic for $p>3$.

Set $C=Z_1/8$. From the raw moments divided by $8$, compute the central
moments $u_2,u_4,u_6$ about $C$ by binomial expansion. Put
\[
 s_1=u_2,\qquad s_2=(3s_1^2-u_4)/2,\qquad
 s_3=(u_6-15u_4u_2+30u_2^3)/16 .
\]
Expanding the even powers and cancelling odd sign monomials gives
$s_j=\sum_i a_i^{2j}$. The elementary symmetric functions of the three
squares are $e_1=s_1$, $e_2=(s_1^2-s_2)/2$ and
$e_3=(s_1^3-3s_1s_2+2s_3)/6$. The octic
\[
 K(z)=\big((z^2-e_1)^2-4e_2\big)^2-64z^2e_3
      =\prod_{\epsilon\in\{\pm1\}^3}\left(z+\sum_i a_i\epsilon_i\right)
\]
follows by pairing the even and odd sign patterns. Their quartics have
the two terms $\pm8a_1a_2a_3 z$ and the same even part.
Thus $K(C)$ is the desired product. Repeated division by $z-C$ gives
the exact number of zero determinants even in small odd characteristic.
This readout works across curves and targets without training.

The verifier checks both assertions on all 146,784 eligible such seed
triples of its order-19 experiment, including 75,840 with a hit. Plugging
pooled moments into the formula on its other nonempty regular inputs
gives 171,888 errors in 220,320 cases. Multiple abscissa choices in
one slot mix different affine sign sums; their pooled central moments
do not recover the three squares above. At the tuned large widths the
chance even one slot has at most two points is $o(\lambda^2)$ by the
slot variance bound. The single-orbit formula supplies no running-time
improvement at the required arity.

\section{What the exact finite distributions show}

We enumerate fresh seeds with their original multiplicities.
For a fixed target and prescribed observation, group the regular seeds
by observed values. The exact minimum number of mistakes is the sum,
over cells, of the smaller of the positive and negative seed masses.
This is an information measurement with a fully labeled table.
The following rows give it the individual slot sizes as side information.
The baseline column uses those sizes alone. The last two columns use
the first six moments. All rows concern existence.

\begin{center}
\begin{tabular}{rrrrrrr}
\toprule
$p$ & $m$ & width & Regular seeds & Baseline & Zero flags & Full values\\
\midrule
13&2&2&10,736&2,208&1,728&0\\
23&2&4&157,572&38,760&37,808&0\\
13&3&2&884,864&108,816&84,192&0\\
\bottomrule
\end{tabular}
\end{center}

A larger experiment uses $E:y^2=x^3+4x+1$ over $\F_{61}$, order 67,
$m=2$, width 2 and target $T$ at cycle index 1 in the verifier. The test
population consists of all regular seeds with $D>6$. Of the 1,180,368
qualifying seed pairs, 171,272 admit a relation. The observation
$(D,Z_1,\ldots,Z_6)$ has 70,788 distinct values, each with a single
existence label. The always-zero error is 14.51 percent here.

A dictionary learned on every other nonzero target except $\pm T$,
including both signs in the raw training set, errs on 171,168 test
seeds when an unseen key defaults to zero. Normalizing $Z_j$ by $y_T^j$
gives 171,064 errors. Training on other targets therefore recovers very
little on this test. A target-specific dictionary can instead be trained
by direct planting. Keep the observed positive signatures from eligible
plants and predict zero on unseen signatures. Its expected error on the
test population, over 1,000, 10,000 and 100,000 independent
\emph{unconditioned} plants at $T$, is 12.2666, 4.7812 and 0.9555
percent respectively. These numbers use the exact occupancy formula
$\sum_{v:\,Y(v)=1}P(v\mid{\rm test})(1-Q(v,{\rm test}))^{n_Q}$.
Purity was verified by enumeration for this fixture. A deployed lookup
only needs the planted signatures and its membership tests.
A reproducible 100,000-plant trial trained this way stored 12,377
positive signatures. It missed 11,172 of the 171,272 positive test
seeds and made zero false-positive predictions, an overall error of
0.9465 percent on the specified 1,180,368 seeds. The verifier caches
feature evaluations on the toy mask family. Its training loop receives
the moment signatures, the eligibility bit and the planted draws.

We also measured a cheaper readout by making fixed random affine
projections of the six field values, retaining $D$. On the same test
population the exact counts below compare the best decision on each bin
with the decision learned from the \emph{population} planted score.
Projection coefficients are fixed before observing labels.
\begin{center}
\begin{tabular}{lrrr}
\toprule
Observation & Occupied bins & Optimal mistakes & Planted-score mistakes\\
\midrule
One field residue & 122 & 171,272 & 171,272\\
Two field residues & 7,369 & 165,028 & 165,760\\
Three field residues & 63,923 & 14,816 & 14,844\\
Ten quadratic-residue bits & 2,048 & 171,128 & 171,180\\
\bottomrule
\end{tabular}
\end{center}
Only this set of projections was tested. Three retain substantial
information but already distinguish nearly as many cells as the full
six-moment observation. These finite rates concern the stipulated subset
of the seed law and do not establish the growing-arity guarantee.

\paragraph{A test on larger fields at balanced widths.}
With three slots, take the order-103 curve $y^2=x^3+17x+1$ over $\F_{101}$
and the order-1019 curve $y^2=x^3+9x+1$ over $\F_{1031}$.
Widths $(4,4,4)$ and $(8,16,16)$ give $\lambda=0.3238$ and $0.2476$,
respectively, both in the specified dyadic range. At cycle target 1 on
each curve, draw 100,000 plants, 25,000 independent unlabeled training
seeds and 20,000 further fresh test seeds. Train the six-moment plus $D$
dictionary from eligible plants. For compressed features, fit the
empirical planted score; all methods predict no hit on irregular inputs.
The table gives mistakes on the 20,000 test seeds (including the irregular
ones). Counts here are samples, unlike the exhaustive results above.
\begin{center}
\begin{tabular}{lrr}
\toprule
Readout & $\F_{101}$ & $\F_{1031}$\\
\midrule
Always absent & 4,668 & 4,068\\
Six-moment dictionary and $D$ & 4,653 & 4,068\\
Tuple count $D$ alone & 4,063 & 4,013\\
One field projection and $D$ & 4,246 & 6,593\\
Two field projections and $D$ & 5,648 & 4,081\\
\bottomrule
\end{tabular}
\end{center}
The two dictionaries stored 67,004 and 89,137 distinct planted
signatures. The regular-only dictionary errors were 3,211 of 15,618
and 3,640 of 18,465, against 3,226 and 3,640 true positives,
respectively. At this training budget, the larger test shows no gain
over predicting absence. The probe also checks each regular test label
against direct group addition. These are fixed-target experiments. Six
moments were chosen for all the finite comparisons. The asymptotic
interpolation bound uses the growing prefix specified in the first section.

\section{The true unweighted spectrum on fresh seeds}

Fix $T$, let $n_v$ be the number of tuple determinants with value $v$,
and set $C=N\lambda=\prod_i\nu_i$, $A_0=\prod_i(1+2/\nu_i)$.
Let $R_T$ denote regularity.
\begin{proposition}
The expected number of ordered, unequal tuple pairs with the same field
value, restricted to regular seeds, satisfies
\[
 \E\left[1_{R_T}\sum_v n_v(n_v-1)\right]
 \le d A_0 C^2/N+md(2^m-1)C/N .
\]
For two fixed eligible targets $S\ne\pm T$ on the same masks, their
histograms satisfy
\[
 \E\left[1_{R_T\cap R_S}\sum_v n_v(T)n_v(S)\right]
 \le d A_0 C^2/N+md\,2^m C/N .
\]
\end{proposition}
\begin{proof}
Classify tuple pairs by the exact set $J$ of slots with matching
abscissas. If $J$ is proper, choose the first tuple, the signs on $J$,
and every remaining point of the second tuple except one. Regularity
ensures the $m-1$ fixed points of that tuple and its fixed target column
are distinct. They impose $m$ independent conditions on $L(dO)$ since
the remaining divisor degree is one. As a function of the last point,
the determinant is therefore a nonzero section of $L(dO)$. It vanishes
at the fixed points, hence is nonconstant. Any specified field value has
at most $d$ preimages. There are at most
$d\,2^{|J|}N^{2m-|J|-1}$ such pairs. The two-abscissa hash law
bounds joint acceptance by $\prod_{i\in J}\rho_i\prod_{i\notin J}\rho_i^2$.
Summing gives the first term in each inequality.

If all abscissas match, the points of the second tuple are fixed sign
flips of the first. In the same-target inequality omit the identical
tuple. For any other sign pattern the section
$\Delta_{\epsilon\mathbf a}(T)-\Delta_{\mathbf a}(T)$, of degree at
most $d$ separately in each curve coordinate, is nonzero. Set one
flipped $a_i=-T$. One determinant vanishes and the other contains
distinct columns $T,-T$. They may be completed to a nonsingular
evaluation matrix by varying the remaining points, since $d\ge3$ and
the affine evaluations span (already on the rational curve for $N-1>d$).
For two different targets $S\ne\pm T$, set any $a_i=-T$. Its signed
image is distinct from $-S$ for every sign pattern, giving the same
argument even for the identity pattern.

Such a nonzero multisection has at most $md N^{m-1}$ affine rational
zeros. To see this, fix a rational tuple where it is nonzero. Specialize
one coordinate to its value there. By induction at most
$(m-1)d N^{m-2}$ choices of the other coordinates kill the specialized
section. At each other choice the remaining one-variable function has
at most $d$ zeros. This proves the stated bound, also allowing constant
slices. All-matched tuples are kept with probability $\prod_i\rho_i$.
Summing over sign patterns finishes the proof.
\end{proof}

Distinct slot abscissas have nonpositive indicator covariance, so
$\operatorname{Var}|F_i|\le2\E|F_i|$. The slots are independent and
$\E|F_i|=(1-1/N)\nu_i$. Thus $D/C$ tends to one in probability, with
failure probability smaller than any fixed power of $\lambda$ for any
fixed relative tolerance. In particular $D<p$ with that probability.
Since $d\lambda=o(1)$, the proposition gives $o(C)$ expected repeated
occurrences on regular roots. The number of distinct nonzero values is
$(1-o_{\Prb}(1))C$ on fresh seeds. Likewise the two spectra at $S\ne\pm T$
have $o_{\Prb}(C)$ values in common. Consequently subtracting the
moments of one fixed reference target generally leaves a long sequence
to reconstruct.

On true-positive seeds the following weaker bound suffices. Write
$J_*= \sum_v n_v(n_v-1)$. The unconditional collision estimate
restricted to regular seeds is $O(dC\lambda)$. Dividing by
$P(Y=1)=\Theta(\lambda)$ gives conditional expectation $O(d C)$ there.
Apart from conditional probability at most $\eta+o(1)$, the seeds are
regular, $D\asymp C$, $D<p$, and $J_*=O(dC/\eta)$.
Cauchy--Schwarz then gives at least $\Omega(\eta C/d)-1$ distinct
nonzero values. For instance $\eta=1/L$ gives $2^{b-O(\log L)}$ values
on a $1-o(1)$ fraction of the positives. The stronger unconditional
near-injectivity estimate has not been asserted after this conditioning.

\paragraph{The cost of extrapolation.}
When $D<p$, every represented frequency survives in the field.
The infinite sequence $Z_j$ for $j\ge1$ has minimal linear recurrence
order precisely the number of distinct nonzero values. Indeed a
recurrence with characteristic polynomial $q$ gives
$\sum_{v\ne0} n_v v^j q(v)=0$ for every $j\ge1$. The first as many
equations as distinct nonzero values form an invertible Vandermonde
system after nonzero column scaling. Thus $q(v)=0$ at every represented
nonzero $v$, and their product of linear factors attains the degree.
Extrapolating via an explicit minimal recurrence or reconstructing the
whole spectrum by Prony interpolation therefore exceeds the proposed
budget even on positive roots. A decoder of the single zero count need
not construct such a recurrence.

\section{Scope of the distributional obstructions}

\paragraph{Random weights.}
One proposed alternative draws independent uniform field weights for
each signed point, slot by slot, and reuses them for all $k$ moments.
Use a product of the chosen weights in each tuple term. Fix typical
regular masks and one point in every slot except the last. The
determinant of the resulting slice is a nonconstant function of degree
at most $d$, so it takes at least $k+1$ distinct values if
$|F_m|>dk$. This size condition fails with probability $o(\lambda^2)$.
Conditional on all other weights, the vector of weighted moments
of orders $0,\ldots,k$ is a linear map of the last slot's random
weights. A $(k+1)$-rowed minor of this map is a nonzero polynomial in
the other slots' weights. In fact specializing those weights to the
chosen points produces a Vandermonde minor on $k+1$ distinct values.
The minor has total degree at most $(k+1)(m-1)$. Hence, except with
probability at most $(k+1)(m-1)/p$, the linear map is onto and the full
field vector is exactly uniform. The conclusion concerns the statistic
when the final weights are hidden. It gives no lower bound for unit
weights or for algorithms making joint use of all realized weights.

\paragraph{What near-injectivity alone implies.}
As a control suppose the spectrum were a uniform $t$-subset of $\F_p$,
$t/p=\pi\asymp1/L$. Its first $k$ field moments carry at most $k\log_2 p$
bits. Translating the underlying set by $-x$ transforms its moments by
the binomial theorem. Any decoder of membership at zero with error
$\varepsilon$ therefore predicts membership at each $x$ with that same
error under this control distribution. The entropy chain rule and
concavity give
\[
 \log_2 {p\choose t} \le k\log_2p+p h_2(\varepsilon).
\]
Using $\log_2{p\choose t}\ge ph_2(\pi)-\log_2(p+1)$ and concavity of
binary entropy, for $\varepsilon\le\pi$ we obtain
\[
 \varepsilon\ge \pi-
 \frac{k\log_2p+\log_2(p+1)}
      {p\log_2((1-\pi)/\pi)} .
\]
This applies even to arbitrary nonlinear moment readouts of a random
subset. The actual masks contain only $O(mb)$ seed bits and do not have
this translation law. The finite decoding results illustrate why that
distinction matters.

\paragraph{End-to-end comparison.}
The desired fresh evaluator has error $O(L^{-2})$ and logarithmic
cost $O(mb/(\sqrt L\,g))=O(b/m)$ at the proposed arity.
The planted score makes learning and testing short observations compatible
with that cost. The finite lookup rates and the recurrence lower bounds
do not prove an evaluator meeting it. The best complete project estimate
therefore remains
\[
 \log_2\E T_{\rm solve}\le
 (26.2933775156+o(1))\,b/L^{1/4}
\]
under packet assumption A.1 and the inherited contraction and reduction
estimates in \texttt{../short\_hash\_count/count.tex}.
Run \texttt{python -B research/joint\_parity/check\_moment\_decoding.py}
from the project root for the exact tests and dictionary trial, and run
\texttt{python -B research/joint\_parity/probe\_planted\_moments.py}
for the independent larger-field training and test draws.
\end{document}
